% Generated by report/make_restated.py from theory/B2_effective_curvature.tex: do not edit (edit the note and rerun).
\begin{restated}{Remark}{rem:B2-gd-scope}{B2}{0}{scope of Proposition~\ref{prop:B2-gd}}
(i) The alignment time is of order $1/\log(1/\rho_{\mathrm{gap}})\approx1/(\gamma(\lambda'-\lmin))$ iterations;
when the spectrum of $H$ is dense near $\lmin$ this can exceed the horizon of an experiment, in which case
$\muEff(\theta_t)$ decreases monotonically towards $\lmin$ without reaching it: monotonicity (b) is the
robust prediction, the limit (c) the asymptotic one. (ii) The case $Pv_0=0$ is not exceptional: by
Lemma~\ref{lem:B2-span}, for logistic regression with $d>n$ and $\theta_0\in S$ the bottom eigenspace $E_{\min}=S^\perp$
of $H_\ast$ is orthogonal to the whole trajectory, and in the quadratic model (d) then gives the alignment with the
flattest direction of $H|_S$, not with $E_{\min}$. Both statements hold in exact arithmetic. Rounding errors re-inject a
component along $E_{\min}$; that this component remains invisible, except possibly close to the precision floor, is
\emph{not proved} here: it is the heuristic of Remark~\ref{rem:B2-gd-rounding}. (iii) For
$\gamma\in(1/\lambda_1,2/\lambda_1)$ the contraction factors are $\abs{1-\gamma\lambda_k}$ and the largest one may
be attained at $\lambda_1$ (for $\gamma=2/(\lambda_1+\lmin)$ both extreme components contract by the same factor
$(\lambda_1-\lmin)/(\lambda_1+\lmin)$, the top one with alternating sign), so (b)--(c)
fail; the project uses $\gamma_{\mathrm{GD}}\le1/L=1/\lambda_1$, as in the book.
\end{restated}
